\section{Six ways to construct a shared seam}
\label{sec:seams}

A cut between two pieces gives each the outline of its neighbor. The six styles
in the custom maker differ in how they construct that line. For the first three,
the starting regions already fill the board before the characteristic interlocks
develop. Figure~\ref{fig:growing-seams} follows a single edge through that change.

\begin{figure}[!ht]
\centering
\begin{minipage}{.98\linewidth}
\textbf{(a) Dendrite.}
Seeds spread until their territories meet. Small waves disturb the new boundary.
The 2012 solver gives adjoining pieces different phase labels and reuses five
labels across the board. Each phase carries a temperature field. As a tip advances,
it releases heat that slows growth along its flanks. That feedback lets branches
develop on either side of the seam. Over long runs, heat can spread between
distant pieces sharing a label and soften the branches~\cite{puzzle2012}.
\par\smallskip
{\centering\plate{dendrite}{.89}\par}
\medskip
\textbf{(b) Amoeba.}
The studio uses the same crystal-growth model with a different parameter set.
The paper identifies seed spacing, ripple wavelength and amplitude, and
supercooling as controls over the resulting shapes. The published description
of Amoeba describes short, closely spaced projections and a higher piece density.
Its numerical preset remains unpublished~\cite{puzzle2012,cutguide}.
\par\smallskip
{\centering\plate{amoeba}{.89}\par}
\medskip
\textbf{(c) Maze.}
Seed regions first expand around the outline and fixed figures. Their shared
edges become elastic rods. Each rod tends to straighten, but its length increases
inside a confined area. The edges buckle, collide and push one another into the
long folds visible in the finished wood~\cite{maze}.
\par\smallskip
{\centering\plate{maze}{.89}\par}
\end{minipage}
\caption{Three routes from initial territory to an interlocking edge. The first two
rows share the studio's phase-field mechanism. The third illustrates the excess
length accommodated by a growing rod. The drawings are conceptual.}
\label{fig:growing-seams}
\end{figure}

\clearpage
\subsection{Curling a rod, joining arcs and trading cells}

Wave adds a force to the elastic edges of Maze. Coral and Matrix instead compare
candidate arrangements and score the interlock they create. The unit of change
in Figure~\ref{fig:constructed-seams} is, respectively, a growing rod, an arc
chain or a grid cell.

\begin{figure}[!ht]
\centering
\begin{minipage}{.98\linewidth}
\textbf{(a) Wave.}
A directional Voronoi layout elongates the initial pieces along a flow field.
Each edge receives vortices according to its length. They move with the edge as
the rod lengthens, drawing it into the spirals and hammerhead-shaped hooks
characteristic of this cut~\cite{wave,cutguide}.
\par\smallskip
{\centering\plate{wave}{.89}\par}
\medskip
\textbf{(b) Coral.}
Jules Khan's generator starts with neighboring cells and proposes replacement
edges made of joined circular arcs. Random walks yield several routes between
the same endpoints. The search scores arc chains for length and interlock,
avoids self-crossings and leaves enough material in narrow parts of a piece
for it to hold together~\cite{coral,cutguide}.
\par\smallskip
{\centering\plate{coral}{.89}\par}
\medskip
\textbf{(c) Matrix.}
The outline, fixed figures and guide curves orient a mesh of mostly quadrilateral
cells. The initial pieces each take roughly the same number of cells. The
optimizer looks several cell swaps ahead as it balances size and interlock. To
score a shared edge, it maps its tangent directions onto a circle. The studio
treats coverage beyond \(180^\circ\) as evidence of interlock. A range approaching
\(360^\circ\) indicates a stronger grip~\cite{matrix,cutguide}.
\par\smallskip
{\centering\plate{matrix}{.89}\par}
\end{minipage}
\caption{Three constructions at edge scale. The red marks in (a) represent
vortices. In (b), gray arcs trace alternative routes and blue marks the selected
seam. In (c), the highlighted cells change ownership across the boundary.}
\label{fig:constructed-seams}
\end{figure}
